\documentclass[12pt]{article}
\usepackage[utf8]{inputenc}
\usepackage[T1]{fontenc}
\usepackage{hyperref}
\title{Tampuan Language Resources for K-12 Curriculum Development}
\author{Compiled by Claude Sonnet 5 (claude-sonnet-5) for Andr\'as Kornai}
\date{}
\begin{document}
\maketitle

\section{Introduction}
Tampuan (ISO 639-3 tpu; Glottocode tamp1251) is an Austroasiatic language of Cambodia. Speaker estimates in the sources I found range from about 31,000 (2008 census) to 57,000 (2013 survey). The search found a Tampuan-English dictionary and some documentation of a Khmer-based script. It found no Tampuan textbooks or primers that I could link to directly.

\section{Primary Educational Resources}
The sources describe bilingual education programmes for Tampuan run by ICC, UNESCO and CARE, using a modified Khmer script that MOEYS approved in 2003. I found no downloadable primers or textbooks from these programmes. One paper on Tampuan education says teaching and learning materials are scarce.
\begin{itemize}
\item Walter Dawson's paper on Tampuan education (PDF): \url{https://icu.repo.nii.ac.jp/record/1403/files/KJ00004933597.pdf}
\end{itemize}

\section{Linguistic and Reference Materials}
\begin{itemize}
\item Online Tampuan-English dictionary on Webonary, with just over 2,000 entries at the time of the snapshot: \url{https://www.webonary.work/tampuan/overview/introduction/}
\item A further Webonary entry page: \url{https://www.webonary.work/tampuan/?p=9444}
\item The printed \emph{Tampuan Khmer English Dictionary} (2007) by Crowley, Vay Tieng and Wain Churk, held by the Center for Khmer Studies Library: \url{https://library.khmerstudies.org/bib/1674}
\item Tampuan numerals (Max Planck Institute for Evolutionary Anthropology): \url{https://lingweb.eva.mpg.de/channumerals/Tampuan.htm}
\item Wikipedia overview of the language: \url{https://en.wikipedia.org/wiki/Tampuan_language}
\item Wikipedia overview of the people: \url{https://en.wikipedia.org/wiki/Tampuan_people}
\item Lexical record in the Huma-Num (CRDO) archive: \url{https://cocoon.huma-num.fr/exist/crdo/lexvo.xql?code=tpu}
\end{itemize}

\section{Religious and Community Resources}
The search did not surface any Tampuan scripture or community-organisation pages. The organisations named in the sources are EMU International, International Cooperation for Cambodia (ICC), CARE and MOEYS, but I found no web pages for them.

\section{Digital and Audio Resources}
The only digital resource found is the Webonary dictionary above. I found no audio or video resources.
\begin{itemize}
\item Omniglot page on the Tampuan writing system: \url{https://omniglot.com/writing/tampuan.htm}
\end{itemize}

\section{Curriculum Development Context}
Tampuan has an approved Khmer-based orthography (MOEYS, 2003), a modest dictionary and earlier bilingual-education pilots. Published curriculum materials appear to be scarce, so a K-12 curriculum would need new material. The first steps would be to contact the ICC, CARE and EMU programmes about their existing materials and to build on the Webonary dictionary.
\end{document}